\section{Conclusion}
\label{sec:conclusion}

We presented a two-stage deep learning pipeline for the automatic detection
and subtype classification of intracranial hemorrhage in non-contrast head
CT scans, reproducing and extending the first-place RSNA 2019 solution on
the publicly available CQ500 dataset.
Stage~1 employs an ensemble of three heterogeneous CNN backbones
(DenseNet-121, DenseNet-169, SE-ResNeXt-101) that classify each CT slice
using a three-channel adjacent-slice representation produced from CT
windowing.
Stage~2 refines the per-slice ensemble logits using a bidirectional GRU
sequence model with a parallel dilated convolutional branch, additionally
incorporating slice thickness as a DICOM metadata feature---a detail
described in the original paper but commonly omitted in reproductions.
Calibrated uncertainty estimates (95\,\% bootstrap confidence intervals)
are reported for all key metrics.

Experimental results on the CQ500 test split show that the ensemble
outperforms any individual backbone, and that the sequence model improves
macro-AUC by a further margin, consistent with the gains reported by the
original authors.
Including slice thickness provides an additional incremental benefit,
particularly for rare subtypes with high inter-slice variability.

Several directions offer promising avenues for future work.
First, the current system propagates study-level labels to all slices,
introducing label noise on non-annotated slices of positive studies;
a multiple-instance learning formulation could address this and potentially
improve slice-level calibration.
Second, the three-backbone ensemble could be extended with more diverse
architectures---such as vision transformers or hybrid CNN–transformer
models---which may provide further complementary representations.
Third, the REST API and Docker deployment developed alongside this work
provide a foundation for prospective clinical evaluation, in which the
system would be integrated into a DICOM viewer to flag high-priority studies
in the radiology worklist.
Finally, the pipeline could be adapted to other cranial CT findings (e.g.,
mass effect, midline shift) with minimal modification, owing to its modular
DVC-managed data pipeline and configurable params.yaml interface.
